\section{Query access and classical output}\label{query:section}
A formulation can have few cone factors, a simple Newton system, and an
inexpensive normalized quantum state, while its requested classical output
still requires reading most of the input. We show this on the instances
used for the geometric results. The query tools are standard: the
polynomial method, search, approximate counting, and quantum information
bounds \citep{Beals2001,BBBV1997,NayakWu1999,BHMT2002,Nayak1999}. What is
new is their application to specified conic lifts, barriers, and output
requirements; we do not claim new Boolean query bounds.

\subsection{Access conventions and an information bound}
The input is a string $b\in\{0,1\}^n$, with $\sigma_j=(-1)^{b_j}$,
accessed through the clean indexed oracle
\[
 O_b|j,z\rangle=|j,z\mathbin\oplus b_j\rangle.
\]
Quantum and randomized query complexities below have worst-case success
probability at least $2/3$. \emph{Public} means independent of $b$.
Public arithmetic, magnitude preparation, index routing, numerical
precision, and writing the output are counted separately. An explicit classical vector means a list of real coordinates at the
specified accuracy, not a circuit, a sampling data structure, or
an amplitude state. For a nonzero real vector $v$, its amplitude state is
$|v\rangle=\norm v^{-1}\sum_jv_j|j\rangle$ in an ordinary complex
Hilbert space.

Sample-and-query access supplies coordinate queries, squared-coordinate samples, and the specified global norm.
Coherent preparation access is a fixed unitary completion
$U_b=O_\sigma U_{\rm pub}$ of the preparation, where $U_{\rm pub}$ is
public and $O_\sigma$ is diagonal: it multiplies specified basis states
(slots) by their hidden signs and fixes all other slots. Both $U_b$ and its
inverse cost a constant number of queries to $O_b$.
A reference slot on which the oracle acts trivially may be added, so
controlled phase queries and clean bit queries are equivalent up to a
constant factor. The model excludes free norms of hidden-dependent
projections, free postselection probabilities, and arbitrary
hidden-dependent unitary completions.

\begin{lemma}[Explicit approximate sign recovery]\label{query:recovery}
Suppose a $T$-query algorithm, on each $\sigma\in\{-1,1\}^n$, returns a
classical string with at most $\alpha n$ errors with probability at least
$2/3$, where $0\leq\alpha<1/4$ is fixed. Then $T=\Omega_\alpha(n)$.
\end{lemma}
\begin{proof}
Purify the algorithm and defer its measurements. Induction on its clean
bit queries gives
\[
 |\psi_\sigma\rangle=
 \sum_{S\subseteq[n],\ |S|\leq T}\left(\prod_{j\in S}\sigma_j\right)|a_S\rangle.
\]
Indeed a bit query, in its target's $|+\rangle,|-\rangle$ basis, either
preserves a coefficient or multiplies it by one $\sigma_j$; input-independent
unitaries preserve degree. Thus the ensemble of final states is supported
on a subspace of dimension $D_T\leq\sum_{j=0}^T\binom nj$. For uniform
$\sigma$, the expected fraction of erroneous decoded bits is at most
$d=(1+2\alpha)/3<1/2$. If $\widehat\sigma$ is the decoded string, elementary
conditional-entropy bounds and concavity of the binary entropy $H_2$ give
\[
 H(\sigma\mid\widehat\sigma)
 \leq\sum_j H(\sigma_j\mid\widehat\sigma_j)\leq nH_2(d).
\]
Holevo's information bound therefore implies
\[
 n[1-H_2(d)]\leq I(\sigma;\widehat\sigma)\leq\log_2D_T.
\]
For $T\leq n/2$, the right-hand side is at most $nH_2(T/n)$, proving a
linear lower bound; larger $T$ is trivial. These are standard entropy
arguments underlying quantum random access bounds
\citep{Nayak1999}; the Fourier span also appears in oracle interrogation
\citep{vanDam1998}.
\end{proof}

We also use the standard bounded-error complexities of estimating
the mean of $n$ bits to additive error $\delta$:
\begin{equation}\label{query:counting}
 Q(n,\delta)=\Theta\bigl(\min\{n,\delta^{-1}\}\bigr),\qquad
 R(n,\delta)=\Theta\bigl(\min\{n,\delta^{-2}\}\bigr)
 \quad(0<\delta\leq c_0),
\end{equation}
for a sufficiently small universal $c_0>0$. Amplitude estimation, or
reading all bits, gives the quantum upper bound; Nayak and Wu's bound for
separated central weights gives the lower bound
\citep{NayakWu1999,BHMT2002}. For the randomized lower bound, use Bernoulli
inputs of means $1/2\pm C\delta$ when $\delta\geq C'/\sqrt n$.
By concentration, an accurate mean estimate distinguishes the two
distributions, while the divergence between $T$-query transcripts is
$O(T\delta^2)$. This gives $\Omega(\delta^{-2})$. Applying this at
$\delta=C'/\sqrt n$ and monotonicity in $\delta$ give $\Omega(n)$
at finer accuracy; the constants absorb finitely many small $n$. At error strictly below half the
gap between adjacent means, parity gives a direct linear lower
bound \citep{Beals2001}; under the promise of zero or one marked bit,
search instead gives $\Theta(\sqrt n)$ quantum and $\Theta(n)$
randomized queries \citep{BBBV1997,BHMT2002}.

\subsection{Certificate rank, movement, and sign output on one lift}
Fix a simple Euclidean Jordan algebra of rank $r\geq2$ attaining the
dictionary capacity
$B=a(r-1)>0$. Let $n=qB$, $s=n+1$, $q\geq2$, and use full half-Peirce
groups in the lift \eqref{eq:rank-attaining-lift}. Put
\[
 g_\sigma=n^{-1/2}(\sigma_1,\ldots,\sigma_n),\qquad
 v_\sigma=(g_\sigma,0),\qquad X_G^0=c+q^{-1}d.
\]
The reference point is strictly feasible; it need not be the analytic
center.

\begin{theorem}[One instance, different costs for different outputs]
\label{query:eja}
For every hidden string, the minimum total Jordan rank of a genuine
certificate for $v_\sigma$ on this lift is $q$. In the restricted standard
Jordan log-determinant metric,
\begin{equation}\label{query:eja-distance}
 d_F\bigl(X^0,\{X:1-\ip{v_\sigma}{\pi X}\leq\epsilon\}\bigr)
 \geq\sqrt q\log(1/\epsilon),\qquad 0<\epsilon<1.
\end{equation}
One sign query prepares the amplitude state of the projected optimizer.
It also prepares the real-coordinate amplitude states of the structured
lifted optimizer and of the standard-barrier central point for every
positive multiplier, provided their public magnitude states can be
prepared. For every fixed $0<\epsilon<1/8$, outputting an explicit
feasible $\epsilon$-optimal point, projected or lifted, requires
$\Theta_\epsilon(n)$ quantum and randomized queries.
\end{theorem}
\begin{proof}
Every group has $\norm{g_G}^2=1/q$. The unique projected optimizer has
$w_G=g_G$, $b=0$, and feasibility forces $z_G=1/q$. Set
\[
 A=\frac1{2q},\qquad
 Y_G=Ac-\frac{g_G}{\sqrt2}+\frac{P(g_G)c}{2A}.
\]
Lemma~\ref{lem:perspective} gives a positive rank-one element with tail
trace $1/2$. Direct pairing gives
$\sum_G\ip{Y_G}{X_G}=1-g_\sigma^Tw$ on the whole affine slice.
Every certificate has nonzero half-Peirce coefficient in each full group,
so every block is nonzero, proving the exact minimum rank $q$.
The positive eigenvalue of $Y_G$ is $(q+1)/(2q)$. Its primitive support
idempotent $u_G$ has $c$-mass $1/(q+1)$ and $d$-mass $q/(q+1)$, hence
\[
 \det_{u_G}Y_G=\frac{q+1}{2q},\qquad
 \det_{u_G}(P(u_G)X_G^0)=\frac2{q+1}.
\]
Thus the scale in Theorem~\ref{thm:movement-minor} is
$\Delta=q[(1/q)^q]^{1/q}=1$, proving \eqref{query:eja-distance}.

The optimizer state is $n^{-1/2}\sum_j\sigma_j|j\rangle$.
For central points, the determinant formula
\eqref{eq:perspective-det} and stationarity in $w_G$ give
$w_G=t_Gg_G$ with $t_G>0$. The objective and barrier depend on group
orientations only through these inner products and norms. Equal group
norms and strict convexity imply common public values $t_G=t$ and $z_G=z$,
and a public $b$. All remaining coordinates are public. Phase kickback
therefore supplies the hidden signs in one query after public magnitude
preparation. For the Albert algebra, this statement encodes its 27 real
coordinates; it assumes neither octonionic quantum amplitudes nor an
efficient Jordan-product circuit.

For any feasible projected point $x$ of support gap at most $\epsilon$,
\[
 \norm{x-v_\sigma}^2\leq2(1-v_\sigma^Tx)\leq2\epsilon.
\]
Each incorrectly decoded sign costs at least $1/n$ in this squared error,
so at most $2\epsilon n$ signs are wrong. Lemma~\ref{query:recovery}
applies with $\alpha=2\epsilon<1/4$. Reading all signs constructs the
optimizer and its structured lift, proving the upper bound. Projection
extends the lower bound to explicit lifted output.
\end{proof}

The movement theorem implies at least
$\sqrt q\log(1/\epsilon)/[-\log(1-R)]$ feasible chords of starting local
norm at most $0<R<1$. This bound and the query lower bound hold
simultaneously but do not multiply: the query bound allows reuse of all
information learned at earlier points.

To see what the quantifier reversal \eqref{eq:quantifier-reversal} costs
in the query model, let $u=(n^{-1/2}\mathbf1,0)$, choose a public orthogonal
$U_0$ sending $u$ to the pole, and set
\[
 U_\sigma=U_0(D_\sigma\oplus1).
\]
Then $U_\sigma v_\sigma$ is the rank-one pole support. One coherent use of
this objective-dependent rotation costs one sign query plus public gates.
An explicit classical matrix $\widetilde U_\sigma$ approximating
$U_\sigma$ with operator-norm error strictly below $1/2$ reveals every
sign through the diagonal of $U_0^T\widetilde U_\sigma$, so producing it
costs $\Theta(n)$ queries. Thus the cost of a reformulation depends on
whether it is used coherently or written out; the geometric quantifier
reversal alone does not determine it.

\subsection{Well-conditioned disk checkpoints with hard scalar readout}
For $N\geq2$, minimize $\sum_i c_i^Tz_i$ over $(\B_2^2)^N$, where
$c_i=e_1$ if $b_i=0$ and $c_i=-e_2$ if $b_i=1$. For
\[
 F(z)=-\sum_i\log(1-\norm{z_i}^2),\qquad
 r(a)=\frac{a}{\sqrt{1+a^2}+1},
\]
the central point is $z_i(\eta)=-r(\eta)c_i$. The
Hessian of one block is
\begin{equation}\label{query:ball-hessian}
 H(z)=\frac2{1-\norm z^2}I+
       \frac4{(1-\norm z^2)^2}zz^T.
\end{equation}
Its radial-to-transverse eigenvalue ratio is $\sqrt{1+\eta^2}$, so
$\kappa(F''(z(2)))=\sqrt5$, independently of $N$. At this checkpoint the
scalar Hessian is diagonal; its structural graph has only disjoint
constant-size components. With $r_0=r(2)$ and Hamming
weight $w=|b|$, the fixed affine observable
\begin{equation}\label{query:disk-readout}
 Nr_0+\sum_i(z_i)_1=r_0w
\end{equation}
requires $\Theta(N)$ quantum and randomized queries to estimate to any
fixed error strictly below $r_0/2$. Under the promise $w\in\{0,1\}$ the
complexities are instead $\Theta(\sqrt N)$ and $\Theta(N)$. These bounds
follow from parity and search; reading all bits gives the first upper
bound.

The optimal value itself is public and equals $-N$; the observable
\eqref{query:disk-readout} reads out the solution, not the value. An explicit feasible point with total objective gap below $1/8$
also needs linearly many queries, since
\[
 \norm{z_i+c_i}^2\leq2(1+c_i^Tz_i)
\]
places each block within distance $1/2$ of its optimizer, and the two
possible block optimizers $-e_1,e_2$ are at distance $\sqrt2$.
In contrast, one clean bit query prepares $|z(\eta)\rangle$ for any
$\eta>0$: query a uniform index into an axis
register and apply the corresponding public sign. The norm
$\sqrt N r(\eta)$ is public. Ordinary sample-and-query access to the raw
objective adds no information: a sample is a uniformly random block, and
one query reveals its axis. Access to norms of arbitrary coordinate subvectors would be stronger; for example, the
first-coordinate part of the objective has norm $\sqrt{N-w}$.

The affine functional in \eqref{query:disk-readout} has Euclidean norm
$\sqrt N$. Normalizing it to norm one changes its required error to
$O(N^{-1/2})$. Likewise, a point within constant local-norm distance of
the checkpoint need not have the fixed scalar accuracy: the functional
has dual local norm $\Theta(\sqrt N)$, so a generic local-norm error
guarantee must be $O(N^{-1/2})$. The lower bound concerns the specified
scalar output.

\subsection{Query cost versus precision with public magnitudes}
Impose $y_i=x_i$ in each disk and use the raw objective pairs
\[
 \widetilde c_i=\frac{\Gamma}{\sqrt{10N}}(1,3\sigma_i),
 \qquad \Gamma>0.
\]
All raw magnitudes and the raw global norm $\Gamma$ are public. In the
orthonormal coordinates $(x_i,y_i)=t_i(1,1)/\sqrt2$, the reduced domain
is $[-1,1]^N$ and the reduced objective coefficients are
\[
 \beta_i=\frac{\Gamma}{\sqrt{5N}}
 \begin{cases}2,&b_i=0,\\-1,&b_i=1.\end{cases}
\]
Hence
\begin{equation}\label{query:slice-value}
 \operatorname{OPT}(b)=-\frac{2\Gamma\sqrt N}{\sqrt5}
                         +\frac{\Gamma}{\sqrt{5N}}w.
\end{equation}

\begin{theorem}[Query cost versus precision on the product domain]\label{query:precision}
For $0<\epsilon\leq c_0\Gamma\sqrt N$, with $c_0>0$ a sufficiently
small universal constant, estimating \eqref{query:slice-value} to additive
error $\epsilon$ has worst-case query complexities
\begin{equation}\label{query:precision-law}
 Q=\Theta\!\left(\min\left\{N,\frac{\Gamma\sqrt N}{\epsilon}\right\}\right),
 \qquad
 R=\Theta\!\left(\min\left\{N,\frac{\Gamma^2N}{\epsilon^2}\right\}\right).
\end{equation}
The same bounds hold for the objective value at the exact central point
of multiplier $\eta_\Gamma=2\sqrt{5N}/\Gamma$. The reduced Hessian has
condition number at most four along the entire central path. One query
prepares the normalized optimizer state. The normalized projected states
of the checkpoint and of the central predictor there have heralded exact
preparation with $O(1)$ expected queries, or constant-error preparation
with $O(1)$ worst-case queries. An explicit feasible point with objective
gap below $\Gamma/\sqrt{5N}$ requires $\Theta(N)$
queries.
\end{theorem}
\begin{proof}
Equation~\eqref{query:slice-value} is an affine rescaling of $w/N$ with
scale $\Gamma\sqrt{N/5}$. Equation~\eqref{query:counting} proves
\eqref{query:precision-law}, with constants absorbed into $c_0$.
The reduced barrier is $-\sum_i\log(1-t_i^2)$ and its central point is
$t_i=-\operatorname{sign}(\beta_i)r(\eta|\beta_i|)$.
At force $a=\eta|\beta_i|$, the scalar Hessian of the reduced barrier is
\[
 h(a)=\sqrt{1+a^2}\bigl(\sqrt{1+a^2}+1\bigr).
\]
For $a>0$, $h(a)/a^2=s/(s-1)$, $s=\sqrt{1+a^2}$, is decreasing;
thus $h(2a)/h(a)\leq4$, and continuity handles $a=0$.
At $\eta_\Gamma$, the forces are four and two. Subtracting the public central
objective value for $b=0$ gives
\begin{equation}\label{query:checkpoint}
 L_b(\eta_\Gamma)+\frac{2\Gamma\sqrt N}{\sqrt5}r(4)
   =\frac{\Gamma}{\sqrt N}\,\Delta_{\rm cp}w,
 \qquad
 \Delta_{\rm cp}=\frac{\sqrt{17}-\sqrt5}{2\sqrt5}>0.
\end{equation}
This is another affine rescaling of $w/N$ by $\Theta(\Gamma\sqrt N)$.

The optimizer coordinates are $-\sigma_i$, so phase kickback gives their
normalized state. Checkpoint magnitudes are the two public constants
$r(4),r(2)$. Query a uniform index, rotate a flag with amplitude proportional
to the desired signed magnitude, uncompute the input bit, and postselect
the flag. The success probability is uniformly bounded below by
$(r(2)/r(4))^2>0$. Repetition gives heralded exact preparation with
constant expected queries; a fixed number of repetitions gives constant
error. The central predictor is
$\dot t_i=-\beta_i/h(\eta|\beta_i|)$, so the same argument at the specified
checkpoint applies to its two nonzero public magnitudes. This concerns
this particular right-hand side, not arbitrary Newton systems. Finally, any
wrong coordinate sign incurs gap at least
$\min_i|\beta_i|=\Gamma/\sqrt{5N}$. Below that threshold every sign is
correct and all bits are recovered; parity supplies the lower bound.
\end{proof}

Squared-coordinate sampling of the raw objective picks a uniform $i$,
then one of its two coordinates with probabilities $1/10,9/10$; a
returned sign costs one query. The canonical coherent completion
applies the sign phase only to the second coordinate of a public
magnitude preparation. The theorem holds under these forms of access,
but not if $\norm\beta^2=\Gamma^2(4N-3w)/(5N)$ is free,
since that reveals the answer.

\paragraph{Moving the signs into the constraints.}
Equivalently, use the public objective
$\Gamma(1,3)/\sqrt{10N}$ in every disk and impose the hidden equalities
$y_i=\sigma_i x_i$. Their constraint matrix $E$ has rows $(-\sigma_i,1)$
with disjoint supports, hence $EE^T=2I$. Row norms, singular values,
column degrees, and squared-sampling distributions are all public. The
hidden orthonormal coordinates $(x_i,y_i)=t_i(1,\sigma_i)/\sqrt2$ give
exactly the preceding coefficients and query bounds. The canonical coherent
matrix completion applies the sign phases to the first-coordinate slots
of a public preparation; each call costs $O(1)$ queries.
A coherent use of the hidden basis costs $O(1)$ queries; writing it out
as a classical matrix costs $\Theta(N)$.
The ambient optimizer pairs are $(-\sigma_i,-1)/\sqrt2$, so one query
after a public coordinate split prepares their amplitude state. At the checkpoint they are $(-r(4),-r(4))/\sqrt2$ or
$(r(2),-r(2))/\sqrt2$, and the heralded preparation just proved applies.
These statements concern the projected variables; they provide neither
a free hidden-vector norm nor an oracle for the full lifted state.

\paragraph{Comparison with a single direct ball.}
Keep the same raw objective and public equalities $y_i=x_i$, but replace
the product of disks by the one ball $\B_2^{2N}$. The reduced domain is
$\B_2^N$, its standard restricted barrier has parameter one, and
\[
 \operatorname{OPT}_{\rm ball}(b)=-\Gamma\phi(w/N),\qquad
 \phi(u)=\sqrt{(4-3u)/5}.
\]
Since $3\sqrt5/20\leq|\phi'(u)|\leq3\sqrt5/10$, the query complexities
of the optimal value become
\begin{equation}\label{query:direct-precision}
 Q=\Theta\bigl(\min\{N,\Gamma/\epsilon\}\bigr),\qquad
 R=\Theta\bigl(\min\{N,\Gamma^2/\epsilon^2\}\bigr)
 \quad(0<\epsilon\leq c_0\Gamma).
\end{equation}
At multiplier $1/\Gamma$, the central value is
$-\Gamma\phi(u)r(\phi(u))=-\Gamma(\sqrt{1+\phi(u)^2}-1)$.
Its derivative in $u$ is $3\Gamma/[10\sqrt{1+\phi(u)^2}]$, bounded
above and below by positive multiples of $\Gamma$, so the same bounds hold.
By \eqref{query:ball-hessian}, the reduced Hessian has condition number
at most $3/\sqrt5$; it is diagonal plus rank one. The
normalized optimizer and central states all equal $|-\beta\rangle$,
prepared by projecting the raw objective state
onto the public equality subspace. The success probability is
$\norm\beta^2/\Gamma^2\in[1/5,4/5]$, so heralded exact preparation uses
$O(1)$ expected queries. An explicit feasible point with gap
below $\Gamma/(8\sqrt5N)$ determines all signs: its squared distance to
$-\beta/\norm\beta$ is at most $2\epsilon/\norm\beta$, whereas every
optimizer coordinate has magnitude at least $1/(2\sqrt N)$.
The same holds for the public-objective hidden-equality encoding.

For this direct ball, $[\log(\Gamma/(2\sqrt5\epsilon))]_+$ is a lower
bound on the barrier distance from zero to the $\epsilon$-accurate
points. Indeed, if $d=-\beta/\norm\beta$ and gap $g=\norm\beta(1-d^Tt)$, then
$1-\norm t^2\leq2g/\norm\beta$, and the barrier gradient has dual norm
at most one. Thus the logarithmic movement bound and the value query bound hold on the same direct-ball instances. The half-power difference
between \eqref{query:precision-law} and \eqref{query:direct-precision}
compares two different feasible bodies; it does not separate two
formulations of one optimization problem. Moreover, the sliced disk
product is a box, so the factor-minimality results for the full disk
product do not apply to it.

\subsection{Reading out several independent search bits}
An intermediate output between one scalar and a dense vector is one bit
for each of several independent search problems. Let $k$ source balls have even dimension $s\geq4$,
and let each have one hidden marked index $j_a\in[s]$. Define
$\sigma_{ai}=1-2\mathbf1_{i=j_a}$ and minimize
$\sum_{a,i}\sigma_{ai}w_{ai}$. Use the grouped Lorentz formulation
\eqref{eq:grouped-lorentz}, globally smooth and optimal under a
cone-dimension cap: for an integer cap
$3\leq d<s+1$, put
$\ell=\lceil s/(d-2)\rceil$, partition each ball into balanced nonempty
groups, and use
\[
 q_{aG}=t_{aG}-\norm{w_{aG}}^2>0,\quad \sum_Gt_{aG}=1,
 \qquad F=-\sum_{a,G}\log q_{aG}.
\]
For $d\geq s+1$, take $\ell=1$ and the direct ball barrier. Stationarity
at multiplier $\eta_0=4\ell/(3\sqrt s)$ gives exactly
\begin{equation}\label{query:search-center}
 q_{aG}=\frac3{4\ell},\quad
 w_{ai}=-\frac{\sigma_{ai}}{2\sqrt s},\quad
 t_{aG}=\frac3{4\ell}+\frac{|G|}{4s}.
\end{equation}
Indeed $w=-\eta_0q\sigma/2$ and $\ell q+s(\eta_0q/2)^2=1$.
For the public first half $H\subset[s]$, the affine readout $p_a=\sqrt s\sum_{i\in H}w_{ai}+|H|/2$ equals
$\mathbf1_{j_a\in H}$.

Returning all $k$ bits $p_a$ with error less than $1/3$ requires
$\Theta(k\sqrt s)$ quantum and $\Theta(ks)$ randomized queries, whereas
one phase query prepares the projected central state. For the quantum lower bound, on one block use the adversary matrix
joining all indices in $H$ to all in its complement. Its norm is $s/2$;
masking by a single query gives norm $\sqrt{s/2}$. For $k$ blocks, sum its
tensor-factor copies. They share a top eigenvector, giving numerator
$ks/2$; a query masks all terms except its own block, giving the same
denominator. The spectral adversary bound
\citep{HoyerLeeSpalek2007} proves the claim. The matching upper bound finds
all $k$ marks among $ks$ positions: knowing the number remaining, exact
amplitude amplification uses $O(\sqrt{ks/j})$ queries at the $j$-remaining
stage, whose sum is $O(k\sqrt s)$ \citep{BHMT2002}. This avoids the
cost of making $k$ separate searches all succeed. For randomized
algorithms under independent uniform
marks, a randomly chosen block, with all other blocks simulated for free,
has expected query cost $T/k$. Deciding whether its mark lies in $H$
requires $\Omega(s)$ expected queries, hence $T=\Omega(ks)$. Scanning gives the upper bound.

The checkpoint is also well conditioned after an explicit public scaling.
Write $z=\sqrt\ell\,\delta t$, $y=\delta w$, and
$A=2\sqrt\ell\,w_G$. Dividing the Hessian quadratic form by $\ell$ gives
\[
 \frac{16}{9}(z-A^Ty)^2+\frac83\norm y^2,
 \qquad \norm A^2=\ell|G|/s<2.
\]
The coupled two-dimensional matrix, multiplied by nine, has determinant
384 and trace below 72, so all blocks have common eigenvalue bounds
$384/(72\cdot9)$ and $72/9$. The condition number is at most $27/2$,
also after restriction to $\sum_Gz_{aG}=0$. The direct-ball case has
condition number $5/3$. This bound concerns the scaled restricted
Hessian, not an arbitrary augmented KKT matrix. Writing a dense classical
vector still takes $\Omega(ks)$ writes, whereas a list of the $k$ marked
indices keeps the search advantage. These different output requirements
explain why the linear arithmetic cost of the Newton solves does not
settle the query comparison.

\subsection{The same conclusions with one cone factor}
The full disk family has the one-factor PSD formulation
\begin{equation}\label{query:one-psd}
 \begin{pmatrix}S&W^T\\W&I_2\end{pmatrix}\succeq0,
 \quad S_{ii}=1,
 \quad D=S-W^TW,
 \quad F=-\log\det D,
\end{equation}
where the columns of $W$ are the disk variables. The cone order is
$N+2$.
For fixed $W$, the fiber center is
$D=\operatorname{diag}(1-\norm{w_i}^2)$ by Hadamard's inequality, so the
marginal barrier is exactly the disk barrier, with slice parameter $N$.
Consequently, the projected central states and readouts above are
unchanged. The statement on inexpensive state preparation does not cover the
dense auxiliary Gram coordinates $S=W^TW+D$.

Starting at $W=0,D=I$, any feasible endpoint of gap at most $\epsilon$
in the full disk problem satisfies
\[
 \det D\leq\prod_i(1-\norm{w_i}^2)
       \leq(2\epsilon/N)^N.
\]
The height estimate of Lemma~\ref{lem:cm-height}, with parameter $N$,
therefore gives
\[
 d_F\geq\sqrt N[\log(N/(2\epsilon))]_+.
\]
For the public-objective hidden-equality version, let $a_i=|\beta_i|$
and $g_i=a_i(1-t_i t_i^*)$. Since $a_i\geq\Gamma/\sqrt{5N}$ and
$1-t_i^2\leq2g_i/a_i$, the corresponding bound is
\begin{equation}\label{query:packed-distance}
 d_F\geq\sqrt N
 \left[\log\frac{\Gamma\sqrt N}{2\sqrt5\epsilon}\right]_+.
\end{equation}
Both bounds allow arbitrary feasible off-diagonal completions along the
path, not just fiber centers. If the initial point is within local
norm $0\leq\rho<1$ of the center, subtract $-\log(1-\rho)$ from the
lower bound before taking the positive part. If each round has at most $m\geq1$ feasible chords ($m$ an integer), each of
starting local norm at most $0<R<1$, dividing by $m[-\log(1-R)]$ gives a
lower bound on the number of rounds.

Alternatively, the disk product is the slice $t=1$ of the single cone
\[
 \mathcal H_{N,2}=\{(t,z):\norm{z_i}\leq t\ (1\leq i\leq N)\}.
\]
By Theorem~\ref{thm:shared-spectral-parameter}, the optimal parameter of
an LHSC barrier on this cone is $N+1$,
attained by
\[
 F_{\mathcal H}(t,z)=-\sum_i\log(t^2-\norm{z_i}^2)+(N-1)\log t.
\]
At $t=1$ its restriction is again the disk barrier. The
public-objective hidden-equality problem is strictly primal feasible at
$t=1,z=0$ and strictly dual feasible: the dual cone consists of
$(M,v)$ with $M\geq\sum_i\norm{v_i}$, so take a sufficiently large positive
$M$ and the public objective blocks as $v_i$.
Therefore, for \emph{any} ambient LHSC barrier,
Theorem~\ref{thm:pd-gap-movement} bounds the number of rounds from an
exact primal--dual central point
of gap $\Delta_0$ to a strictly feasible point of gap
$0<\epsilon<\Delta_0$:
\begin{equation}\label{query:shared-pd}
 T\geq\frac{\sqrt{(N+1)/2}\log(\Delta_0/\epsilon)}
                 {m[-\log(1-R)]}.
\end{equation}
Chords are measured in the primal--dual product metric, and
intermediate points need not be affinely feasible. If the start is at the public multiplier
$\eta_0=2\sqrt5$ and $\epsilon<(N+1)/\eta_0$, the numerator can be
replaced by
$\sqrt{(N+1)/2}\log((N+1)/(\eta_0\epsilon))$.
Indeed $\Delta_0=\nu/\eta_0$, $\nu\geq N+1$, and
$\sqrt\nu\log(\nu/(\eta_0\epsilon))$ is increasing on this range.
The inexpensive-state and diagonal-Hessian statements hold for the
displayed barrier, whereas \eqref{query:shared-pd} holds for every ambient LHSC barrier
with its own metric and central start.

Thus, at $\Gamma=\sqrt N$ and a fixed, sufficiently small error, these
explicit formulations have one cone factor, need $\Theta(N)$ queries for
scalar or explicit output, have projected states preparable with $O(1)$
queries, and need $\Omega(\sqrt N\log N)$ bounded-movement rounds in the
stated metrics. A cost model counting both queries and rounds has their
maximum as a lower bound. Multiplying them would require a separate
theorem showing that every round needs new queries; these fixed-input
families provide no such theorem. This is consistent with quantum
optimization analyses that count data access, linear solving, and
classical output separately \citep{ApersGribling2026}.
